% Worked Examples: Concept 2.6.2 - The Nyquist Criterion
\documentclass[11pt,a4paper]{article}
\usepackage{worked-examples-template}

\title{\textbf{Worked Examples}\\[0.5em]
\large Concept 2.6.2: The Nyquist Criterion\\[0.3em]
\normalsize \coursesubtitle{} -- \coursetitle}
\author{Assoc.\ Prof.\ William Robertson \and Dr Sean McGowan}
\date{\today}

\begin{document}

\maketitle
\tableofcontents
\newpage

\section{Concept 2.6.2: Nyquist Criterion}

\subsection{Example 1: Applying the Nyquist Stability Criterion}

\paragraph{Problem:} Use the Nyquist criterion to determine the stability of a unity feedback system with open-loop transfer function:
\begin{align}
L(s) = \frac{K}{(s+1)(s+2)}
\end{align}

for $K = 5$.

\paragraph{Solution:}

The Nyquist stability criterion states that for a closed-loop system:
\begin{align}
N = P - Z
\end{align}

where:
\begin{itemize}
\item $N$ = number of clockwise encirclements of the point $-1 + j0$ by the Nyquist plot
\item $P$ = number of open-loop poles in the right half-plane
\item $Z$ = number of closed-loop poles in the right half-plane
\end{itemize}

For stability, we need $Z = 0$, which requires $N = P$.

\textbf{Step 1: Count open-loop RHP poles}

The open-loop poles are at $s = -1$ and $s = -2$ (both in left half-plane).
Therefore: $P = 0$

\textbf{Step 2: Analyze the Nyquist plot}

For the frequency response $L(j\omega)$ with $K = 5$:
\begin{align}
L(j\omega) = \frac{5}{(j\omega+1)(j\omega+2)} = \frac{5}{(1+j\omega)(2+j\omega)}
\end{align}

Key points:
\begin{itemize}
\item At $\omega = 0$: $L(j0) = \frac{5}{1 \cdot 2} = 2.5$ (on positive real axis)
\item As $\omega \to \infty$: $L(j\omega) \to 0$ (magnitude approaches 0)
\item Phase is always negative (two poles, no zeros)
\end{itemize}

To check if the plot crosses the negative real axis at $-1$:

At crossover, the imaginary part is zero and real part is negative. For $L(j\omega)$ to be real and negative:
\begin{align}
\text{Im}[L(j\omega)] = 0
\end{align}

Converting to magnitude and phase:
\begin{align}
|L(j\omega)| &= \frac{5}{\sqrt{(1+\omega^2)(4+\omega^2)}} \\
\angle L(j\omega) &= -\arctan(\omega) - \arctan(\omega/2)
\end{align}

The phase crosses $-180°$ when:
\begin{align}
\arctan(\omega) + \arctan(\omega/2) = 180° = \pi
\end{align}

This never happens for finite $\omega$ (only as $\omega \to \infty$, where magnitude is 0).

The Nyquist plot:
\begin{itemize}
\item Starts at 2.5 on the positive real axis ($\omega = 0$)
\item Curves through the lower half-plane (negative imaginary part)
\item Approaches origin as $\omega \to \infty$
\item Never encircles $-1 + j0$
\end{itemize}

Therefore: $N = 0$

\textbf{Step 3: Apply Nyquist criterion}

\begin{align}
Z = P - N = 0 - 0 = 0
\end{align}

\textbf{Result:} The closed-loop system is stable for $K = 5$ since there are no closed-loop poles in the right half-plane ($Z = 0$).
\end{document}
